\clearpage\thispagestyle{otFront}\phantomsection\label{avant-propos}\pdfbookmark[0]{Apprendre et comprendre}{bm-avant-propos}\otPageTitre{Apprendre et comprendre}{Un livre pour toute l’année de Sixième}\fontsize{11}{15}\selectfont

Ce livre accompagne l’entrée au collège : nombres décimaux et fractions,
calcul, proportionnalité, grandeurs, géométrie des distances et des angles,
données, hasard et premiers pas en programmation. Il consolide les acquis
de l’école et prépare le cycle 4, en développant le raisonnement.\par\medskip
\begin{savoirfaire}[Du simple au problème]Chaque chapitre propose ses objectifs,
ses prérequis, deux pages de cours en quatre sections et quinze exercices :
cinq de découverte, cinq d’entraînement et cinq problèmes. Chaque correction
explique la méthode et le sens du résultat.\end{savoirfaire}
\begin{savoirfaire}[Des figures pour comprendre]Demi-droites graduées, bandes,
schémas en barres, rapporteurs, constructions au compas, symétries,
assemblages de cubes, diagrammes et programmes par blocs accompagnent les
notions. Les constructions demandées apparaissent en vert dans les corrigés.\end{savoirfaire}
\begin{savoirfaire}[Chercher et oser]Vingt-et-un défis facultatifs offrent deux
indications séparées et une solution expliquée. Deux bilans mélangent les
chapitres pour apprendre à choisir ses outils.\end{savoirfaire}
Chercher sans réussir du premier coup fait partie des mathématiques :
garder ses essais, les expliquer, puis comparer avec la méthode proposée.

\clearpage
